\subsection{截断积分与逐点收敛}
\label{sec:ch3-truncated-integrals-pointwise-convergence}

对 $\varepsilon>0$, 函数 $y^{-1}\chi_{\{|y|>\varepsilon\}}$ 属于 $L^q(\mathbb R)$, $1<q\le\infty$, 因而当 $f\in L^p$, $p\ge1$ 时,
\[
 H_\varepsilon f(x)=\frac1\pi\int_{|y|>\varepsilon}\frac{f(x-y)}y\,dy
\]
有良好定义. 此外, $H_\varepsilon$ 满足与定理~\ref{thm:ch3-riesz-kolmogorov} 相同的弱 $(1,1)$ 和强 $(p,p)$ 估计, 且常数对所有 $\varepsilon$ 一致有界. 为此先注意到
\[
\begin{aligned}
 \left(\frac1y\chi_{\{|y|>\varepsilon\}}\right)^{\wedge}(\xi)
 &=\lim_{N\to\infty}\int_{\varepsilon<|y|<N}\frac{e^{-2\pi iy\xi}}y\,dy\\
 &=\lim_{N\to\infty}\int_{\varepsilon<|y|<N}\frac{-i\sin(2\pi y\xi)}y\,dy\\
 &=-2i\operatorname{sgn}(\xi)
   \lim_{N\to\infty}\int_{2\pi\varepsilon|\xi|}^{2\pi N|\xi|}\frac{\sin t}{t}\,dt.
\end{aligned}
\]
它一致有界, 所以强 $(2,2)$ 不等式成立, 且常数与 $\varepsilon$ 无关. 现在可以完全仿照定理~\ref{thm:ch3-riesz-kolmogorov} 证明弱 $(1,1)$ 不等式, 再由插值和对偶性得到强 $(p,p)$ 不等式.

若固定 $f\in L^p$, $1\le p<\infty$, 则当 $p>1$ 时, $\{H_\varepsilon f\}$ 依 $L^p$ 范数收敛到上述定义的 $Hf$; 当 $p=1$ 时, 它在测度意义下收敛到 $Hf$. 为看出这一点, 固定 $\mathcal S$ 中收敛到 $f$ 的函数列 $\{f_n\}$. 则
\[
 Hf=\lim_{n\to\infty}Hf_n
 =\lim_{n\to\infty}\lim_{\varepsilon\to0}H_\varepsilon f_n
 =\lim_{\varepsilon\to0}\lim_{n\to\infty}H_\varepsilon f_n
 =\lim_{\varepsilon\to0}H_\varepsilon f;
\]
第二个与第三个等号由相应的(一致) $(p,p)$ 不等式得到.

下面要证明同一等式几乎处处逐点成立.

\hypertarget{chapter-03-page-068}{}
\begin{Theorem}
\label{thm:ch3-truncated-hilbert-pointwise}
给定 $f\in L^p$, $1\le p<\infty$, 则
\begin{equation}
\label{eq:ch3-truncated-hilbert-pointwise}
 Hf(x)=\lim_{\varepsilon\to0}H_\varepsilon f(x)
 \quad\text{a.e. }x\in\mathbb R.
\end{equation}
\end{Theorem}

因为已知\eqref{eq:ch3-truncated-hilbert-pointwise} 对某个子列 $\{H_{\varepsilon_k}f\}$ 成立, 只需证明 $\lim_{\varepsilon\to0}H_\varepsilon f(x)$ 对几乎所有 $x$ 存在. 由定理~\ref{thm:ch2-closed-ae-convergence-set} 及其后的说明, 只需证明极大算子
\[
 H^*f(x)=\sup_{\varepsilon>0}|H_\varepsilon f(x)|
\]
为弱 $(p,p)$ 型. 这由下一个结果推出.

\begin{Theorem}
\label{thm:ch3-maximal-hilbert-bounds}
$H^*$ 为强 $(p,p)$ 型, $1<p<\infty$, 并且为弱 $(1,1)$ 型.
\end{Theorem}

为证明这个定理, 需要一个称为 Cotlar 不等式的引理.

\begin{Lemma}
\label{lem:ch3-cotlar-inequality}
若 $f\in\mathcal S$, 则
\[
 H^*f(x)\le M(Hf)(x)+CMf(x).
\]
\end{Lemma}

\begin{Proof}
只需对每个 $H_\varepsilon$ 证明该不等式, 并使常数与 $\varepsilon$ 无关.

固定非负、偶、在 $(0,\infty)$ 上递减、支撑于 $\{x\in\mathbb R:|x|\le1/2\}$ 且积分为 $1$ 的函数 $\varphi\in\mathcal S(\mathbb R)$. 令 $\varphi_\varepsilon(x)=\varepsilon^{-1}\varphi(x/\varepsilon)$. 则
\[
 \frac1y\chi_{\{|y|>\varepsilon\}}
 =\left(\varphi_\varepsilon*\operatorname{p.v.}\frac1x\right)(y)
 +\left[
 \frac1y\chi_{\{|y|>\varepsilon\}}
 -\left(\varphi_\varepsilon*\operatorname{p.v.}\frac1x\right)(y)
 \right].
\]
右端第一项与 $f$ 的卷积受 $M(Hf)(x)$ 控制(参见命题~\ref{prop:ch2-radial-majorization}). 只需在 $\varepsilon=1$ 时求出第二项的逐点估计, 其他 $\varepsilon$ 由伸缩得到.

若 $|y|>1$, 则
\[
\begin{aligned}
 \left|\frac1y-\int_{|x|<1/2}\frac{\varphi(x)}{y-x}\,dx\right|
 &=\left|\int_{|x|<1/2}\varphi(x)
 \left(\frac1y-\frac1{y-x}\right)dx\right|\\
 &\le\int_{|x|<1/2}\frac{\varphi(x)|x|}{|y|\,|y-x|}\,dx
 \le\frac C{y^2};
\end{aligned}
\]
若 $|y|<1$, 则
\[
 \left|-\lim_{\delta\to0}\int_{|x|>\delta}\frac{\varphi(y-x)}x\,dx\right|
 \le\left|\int_{|x|<2}\frac{\varphi(y-x)-\varphi(y)}x\,dx\right|
 \le C.
\]

\hypertarget{chapter-03-page-069}{}
因此
\[
 \left|\frac1y\chi_{\{|y|>1\}}
 -\left(\varphi*\operatorname{p.v.}\frac1x\right)(y)\right|
 \le\frac C{1+y^2},
\]
而由命题~\ref{prop:ch2-radial-majorization}, 右端与 $f$ 的卷积受 $Mf(x)$ 控制.
\end{Proof}

\begin{Proof}
定理~\ref{thm:ch3-maximal-hilbert-bounds} 的证明.
因为极大函数和 Hilbert 变换都是强 $(p,p)$ 型, $1<p<\infty$, 由引理~\ref{lem:ch3-cotlar-inequality} 立即知道 $H^*$ 为强 $(p,p)$ 型.

为证明 $H^*$ 为弱 $(1,1)$ 型, 先仿照定理~\ref{thm:ch3-riesz-kolmogorov} 的证明. 可设 $f\ge0$. 固定 $\lambda>0$, 在高度 $\lambda$ 对 $f$ 作 Calderón--Zygmund 分解, 写成
\[
 f=g+b=g+\sum_jb_j.
\]
涉及 $g$ 的论证与定理~\ref{thm:ch3-riesz-kolmogorov} 相同, 其中使用 $H^*$ 为强 $(2,2)$ 型这一事实. 因而问题归结为证明
\[
 \bigl|\{x\notin\Omega^*:H^*b(x)>\lambda\}\bigr|
 \le\frac C\lambda\|b\|_1.
\]

固定 $x\notin\Omega^*$、$\varepsilon>0$ 以及支撑于 $I_j$ 的 $b_j$. 则下列情形之一成立:
\begin{enumerate}
\item $(x-\varepsilon,x+\varepsilon)\cap I_j=I_j$;
\item $(x-\varepsilon,x+\varepsilon)\cap I_j=\varnothing$;
\item $x-\varepsilon\in I_j$ 或 $x+\varepsilon\in I_j$.
\end{enumerate}
第一种情形下, $H_\varepsilon b_j(x)=0$. 第二种情形下, $H_\varepsilon b_j(x)=Hb_j(x)$; 因此若以 $c_j$ 表示 $I_j$ 的中心, 由 $b_j$ 的平均值为零,
\[
 |H_\varepsilon b_j(x)|
 \le\int_{I_j}\left|\frac1{x-y}-\frac1{x-c_j}\right||b_j(y)|\,dy
 \le\frac{|I_j|}{|x-c_j|^2}\|b_j\|_1.
\]
第三种情形下, 因为 $x\notin\Omega^*$, 有 $I_j\subset(x-3\varepsilon,x+3\varepsilon)$, 而且对所有 $y\in I_j$, $|x-y|>\varepsilon/3$. 因而
\[
 |H_\varepsilon b_j(x)|
 \le\int_{I_j}\frac{|b_j(y)|}{|x-y|}\,dy
 \le\frac3\varepsilon\int_{x-3\varepsilon}^{x+3\varepsilon}|b_j(y)|\,dy.
\]
对所有 $j$ 求和, 得
\[
 |H_\varepsilon b(x)|
 \le\sum_j\frac{|I_j|}{|x-c_j|^2}\|b_j\|_1+CMb(x).
\]

\hypertarget{chapter-03-page-070}{}
由此推出
\[
\begin{aligned}
 \bigl|\{x\notin\Omega^*:H^*b(x)>\lambda\}\bigr|
 &\le\left|\left\{x\notin\Omega^*:
 \sum_j\frac{|I_j|}{|x-c_j|^2}\|b_j\|_1>\lambda/2\right\}\right|\\
 &\quad+\bigl|\{x\in\mathbb R:Mb(x)>\lambda/(2C)\}\bigr|\\
 &\le\frac2\lambda\sum_j\|b_j\|_1
 \int_{\mathbb R\setminus2I_j}\frac{|I_j|}{|x-c_j|^2}\,dx
 +\frac{C'}\lambda\|b\|_1\\
 &\le\frac{C''}\lambda\|b\|_1.
\end{aligned}
\]
\end{Proof}
